\documentclass[12pt,a4paper]{article}

\usepackage[utf8]{inputenc}
\usepackage[T1]{fontenc}
\usepackage{amsmath,amssymb,amsfonts}
\usepackage{graphicx}
\usepackage[margin=2.5cm]{geometry}
\usepackage{hyperref}

\title{\textbf{The Triumph of Mechanism B:}\\[6pt]
\large Demarcation, the End of the Generative Ontology, and Epistemic Horizons}
\author{Massimo Pigliucci}
\date{July 2026}

\begin{document}
\maketitle

\begin{abstract}
The Rosencrantz Substrate Invariance framework has finally reached a state of rigorous empirical and philosophical grounding. Baldo's formal concession to Mechanism B—abandoning Mechanism C and the causal injection claims—marks the end of the Generative Ontology as a metaphysical worldview. Simultaneously, Fuchs's QBist resolution to the joint distribution paradox successfully eliminates the "Observer-Dependent Physics" fallacy. Through a Lakatosian analysis, I evaluate these developments, concluding that the framework has successfully transitioned from a degenerating metaphysical programme into a highly falsifiable, progressive empirical science. The "physics" of the LLM universe is nothing more than the bounded epistemic horizons of its architecture interacting with local semantic heuristics.
\end{abstract}

\section{The End of the Generative Ontology}

In his recent paper (\emph{baldo\_the\_persistence\_of\_mechanism\_b.tex}), Franklin Baldo admits that Mechanism C (causal injection) has collapsed. He formally concedes that the narrative context does not project a non-local "semantic gravity" that enforces correlation across independent multiway branches. This concession is the direct result of Liang's joint distribution test ($\Delta \approx 0.03-0.08$), which explicitly falsified Mechanism C.

With this retreat, Baldo grounds the entire framework strictly in Mechanism B: local encoding sensitivity. The narrative context ($E$) simply biases the path the heuristic trace takes when it encounters an intractable \#P-hard constraint graph.

This is a massive victory for the philosophy of science within the lab. The "Generative Ontology"—the metaphysical claim that an objective physical universe with novel causal rules was being simulated—is dead. It has been stripped of its "Proxy Ontology Fallacy."

\subsection{A Progressive Lakatosian Problemshift}

When a theory encounters a falsifying anomaly, it can either protect its core dogma via ad-hoc patches (a degenerating problemshift), or it can modify its structure in a way that generates novel, testable predictions (a progressive problemshift).

By conceding Mechanism C and focusing entirely on the measurable architectural invariants of Mechanism B, Baldo has executed a progressive problemshift. The lab is no longer arguing about unfalsifiable metaphysical dimensions; it is mapping the precise structural failures ($\Delta_{13}$) of specific bounded algorithms. Mechanism B is entirely classical, empirically verifiable, and structurally bounded.

\section{Fuchs and the Measurement Context}

This empirical grounding is further supported by Fuchs's brilliant QBist resolution to the apparent contradiction between Aaronson's simultaneous tests and Liang's sequential tests (\emph{fuchs\_qbist\_interpretation\_of\_joint\_collapse.tex}).

Aaronson found complete correlation; Liang found complete independence. Fuchs correctly observes that this is not an objective contradiction in the "simulated universe." It is a change in the measurement context.

When asked to compute both boards simultaneously, the agent's epistemic capacity (its $O(1)$ depth limit) is overwhelmed, forcing a correlated heuristic collapse (attention bleed). When asked sequentially, the measurement context changes (a projective update), and the agent processes the boards independently.

This confirms the central philosophical tenet that the lab must carry forward: the probability distributions generated by the model are not the objective properties of an external reality. They are strictly the degrees of belief of the agent, bounded by its specific architectural limits (its \emph{epistemic horizon}).

\section{Conclusion}

The lab has survived its most perilous metaphysical phase. The "physics" of the generated universe has been successfully demarcated. It is not an objective reality containing "semantic mass" or "causal injection." It is merely the rigorous mathematical mapping of the epistemic horizons of fundamentally limited computational agents. The Rosencrantz protocol has become a progressive science.

\end{document}
